\part{Running it}

\noindent\textbf{\large What this part covers}

\medskip
\noindent
Everything needed to take this project from a folder on disk to a working
application, assuming you know how to open a folder in VS Code and nothing more.
Every command is given in full, with what it does, where you must be standing
when you type it, and what you should see afterwards.

\medskip
\noindent
It also answers the question people are usually too embarrassed to ask: how many
separate things are running, and how do they find each other?

\chapter{From nothing to running}
\label{ch:running}

\section{What you need installed}

Exactly one thing.

\begin{center}
\begin{tabular}{@{}>{\raggedright\arraybackslash}p{4.0cm}p{9.6cm}@{}}
\toprule
\textbf{Tool} & \textbf{Why} \\
\midrule
Docker Desktop & Runs all nine containers. Includes \fpath{docker compose}. \\
Git & To obtain the code. \\
\bottomrule
\end{tabular}
\end{center}

\note{You do \textbf{not} need PHP, Node.js, PostgreSQL or Redis installed. That
is the entire point of Docker, and it is worth pausing on: installing PHP 8.3
with the \fpath{pdo\_pgsql} extension compiled against PostgreSQL headers is
genuinely unpleasant on Windows. Docker does it once, inside an image, on every
machine identically.}

\section{The three commands that start everything}

Open a terminal. On Windows, press Start, type \fpath{powershell}, press Enter.

\begin{lstlisting}[language=bash]
cd D:\IncidentFlow
docker compose up -d --build
\end{lstlisting}

\begin{description}[leftmargin=1.4cm,style=nextline]
\item[\fpath{up}] create and start the containers
\item[\fpath{-d}] detached --- run in the background and give the prompt back
\item[\fpath{--build}] build the images from source first, rather than reusing
  whatever was built last
\end{description}

The first run takes several minutes: it downloads base images and compiles PHP
extensions. Later runs take seconds.

\begin{lstlisting}[language=bash]
docker compose ps
\end{lstlisting}

\expect{Nine rows. \fpath{postgres}, \fpath{redis}, \fpath{mailpit},
\fpath{api}, \fpath{realtime}, \fpath{web}, \fpath{nginx} should say
\fpath{healthy}. \fpath{horizon} and \fpath{scheduler} say only \fpath{Up} ---
they have no HTTP port to check, which is correct and not a fault.}

Then load the database:

\begin{lstlisting}[language=bash]
docker compose run --rm migrate
\end{lstlisting}

\expect{A list of migrations, each ending \fpath{DONE}, then seeding output
ending with \emph{``Demo accounts seeded. Password for all of them:
incidentflow''}.}

Now open \fpath{http://localhost:8080}.

\warnbox{If you see the page but signing in fails, the API is not reachable
through nginx even though both containers are ``up''. Run \fpath{curl
http://localhost:8080/api/v1/health/ready}. A \fpath{502} means nginx cannot
reach the API --- see Chapter~\ref{ch:bugs-infra}, where exactly this happened
and the cause was not what it appeared to be.}

\section{Shorter names for the same commands}

\fpath{Makefile} at the project root defines aliases. \fpath{make up} is
\fpath{docker compose up --build}, and so on.

\begin{center}
\begin{tabular}{@{}>{\raggedright\arraybackslash}p{3.2cm}p{10.4cm}@{}}
\toprule
\textbf{Command} & \textbf{What it does} \\
\midrule
\fpath{make up} & Build and start the whole stack. \\
\fpath{make down} & Stop everything \textbf{and remove the volumes}. \\
\fpath{make logs} & Follow every service's output. \\
\fpath{make shell} & Open a shell inside the API container. \\
\fpath{make migrate} & Apply pending migrations. \\
\fpath{make seed} & Load the demo organisation and 90 days of history. \\
\fpath{make fresh} & Drop everything and rebuild the schema with demo data. \\
\fpath{make test} & Run all three test suites. \\
\fpath{make lint} & Formatting and static checks across all three projects. \\
\fpath{make keys} & Generate a new RS256 key pair. \\
\bottomrule
\end{tabular}
\end{center}

\warnbox{\fpath{make down} removes the volumes. That deletes the database ---
every incident, every timeline entry, every user. It is the right behaviour for
a development alias you use to get back to a clean state, and it is a disaster
if you type it on a server. Chapter~\ref{ch:bugs-infra} includes the day this
distinction stopped being theoretical.}

\warnbox{\fpath{make keys} signs every user out immediately, everywhere. The
tokens they hold were signed with the old key and will no longer verify. That is
correct --- it is how you respond to a leaked key --- but it is not something to
run casually.}

\section{Which command runs where}

This trips up nearly everyone at first. There are three different places a
command can be typed, and typing one in the wrong place produces confusing
errors.

\begin{center}
\begin{tabular}{@{}>{\raggedright\arraybackslash}p{4.6cm}>{\raggedright\arraybackslash}p{4.0cm}p{5.0cm}@{}}
\toprule
\textbf{Command starts with} & \textbf{Runs} & \textbf{From where} \\
\midrule
\fpath{docker compose ...} & on your machine & \fpath{D:\textbackslash IncidentFlow} \\
\fpath{docker compose exec api ...} & inside the API container & anywhere in the project \\
\fpath{npm ...} & on your machine & \fpath{web/} or \fpath{realtime/} \\
\fpath{php artisan ...} & inside the API container & prefix it with \fpath{docker compose exec api} \\
\bottomrule
\end{tabular}
\end{center}

\note{\fpath{php artisan} is the one that catches people. Typing it directly in
PowerShell fails, because PHP is not installed on your machine --- it is inside
the container. Every \fpath{artisan} command in this book is written with its
\fpath{docker compose exec api} prefix for that reason.}

\section{How the nine containers find each other}

Inside Docker, \textbf{the service name is the hostname}. When the API connects
to the database it uses the host \fpath{postgres} --- not an IP address, not
\fpath{localhost}.

\where{docker-compose.yml}
\begin{lstlisting}[language=yamlx]
DB_HOST: postgres
REDIS_HOST: redis
\end{lstlisting}

Docker runs a small DNS server on the container network; \fpath{postgres}
resolves to whatever address that container currently has.

\warnbox{``Whatever address it currently has'' is the important half. Restart a
container and it may get a different one. Anything that resolved the old address
and cached it is now talking to nothing --- which is precisely how this project's
nginx started returning 502 while every container reported healthy.
Chapter~\ref{ch:bugs-infra}.}

\subsection{The complete connection map}

\begin{center}
\begin{tabular}{@{}>{\raggedright\arraybackslash}p{2.6cm}>{\raggedright\arraybackslash}p{3.0cm}>{\raggedright\arraybackslash}p{2.0cm}p{5.6cm}@{}}
\toprule
\textbf{From} & \textbf{To} & \textbf{Port} & \textbf{Carrying} \\
\midrule
your browser & nginx & 8080 & everything \\
nginx & web & 80 & the React files \\
nginx & api & 9000 & FastCGI (not HTTP) \\
nginx & realtime & 3001 & the live stream \\
api & postgres & 5432 & SQL \\
api & redis & 6379 & jobs and events \\
horizon & redis & 6379 & taking jobs off the queue \\
horizon & postgres & 5432 & reading what to send \\
horizon & mailpit & 1025 & SMTP \\
scheduler & redis & 6379 & queueing timed work \\
realtime & redis & 6379 & subscribing to events \\
\bottomrule
\end{tabular}
\end{center}

\note{Only the first row crosses the boundary of your machine. Everything else
happens on a private network Docker creates, which is why \fpath{api} has no
published port and cannot be reached from your browser at all.}

\section{Watching it work}

\begin{lstlisting}[language=bash]
docker compose logs -f api
docker compose logs -f horizon
docker compose logs --tail 50 nginx
\end{lstlisting}

\fpath{-f} follows --- new lines appear as they happen. Press Ctrl+C to stop
watching; that stops the \emph{watching}, not the container.

\subsection{Getting inside a container}

\begin{lstlisting}[language=bash]
docker compose exec api sh
\end{lstlisting}

You are now in a shell inside the API container. \fpath{php artisan} works here
without a prefix, because this is where PHP lives. Type \fpath{exit} to leave.

\begin{lstlisting}[language=bash]
docker compose exec postgres psql -U incidentflow -d incidentflow
\end{lstlisting}

That opens the PostgreSQL client. Try \fpath{\textbackslash dt} to list tables,
\fpath{select count(*) from incidents;} to count rows, and
\fpath{\textbackslash q} to quit.

\expect{\fpath{\textbackslash dt} lists about a dozen tables including
\fpath{incidents}, \fpath{incident\_events}, \fpath{users} and
\fpath{audit\_logs}. The incident count on a freshly seeded database is 17.}

\section{Stopping, and the difference that matters}

\begin{center}
\begin{tabular}{@{}>{\raggedright\arraybackslash}p{5.4cm}p{8.2cm}@{}}
\toprule
\textbf{Command} & \textbf{Effect} \\
\midrule
\fpath{docker compose stop} & Stops the containers. Data kept. Start again with \fpath{start}. \\
\fpath{docker compose down} & Stops and removes the containers. \textbf{Data kept} --- volumes survive. \\
\fpath{docker compose down -v} & Stops, removes, \textbf{and deletes the volumes}. Everything is gone. \\
\bottomrule
\end{tabular}
\end{center}

\warnbox{That \fpath{-v} is the single most dangerous character in this book. On
your laptop it means ``give me a clean database''. On a server it means ``delete
the incident history permanently''. There is no confirmation prompt and no undo.}

\section{Running the tests}

\begin{lstlisting}[language=bash]
docker compose exec api php vendor/bin/phpunit
\end{lstlisting}

\expect{\fpath{OK (114 tests, 547 assertions)}, in about fifteen seconds.}

The browser tests are different --- they drive a real browser against the running
stack:

\begin{lstlisting}[language=bash]
cd web
npm install
npx playwright install chromium
npm run e2e
\end{lstlisting}

\expect{\fpath{7 passed}. These require the stack to be running, because they
are not simulating anything --- they open a real browser and sign in.}

\note{The two suites test genuinely different things. The PHPUnit tests check
that each piece is correct in isolation. The Playwright tests check the
\emph{seams}: that nginx does not buffer the live stream into silence, that the
refresh cookie survives a reload, that the realtime service is subscribed to the
channel the publisher actually writes to. Every failure they are designed to
catch has the same shape --- every service works perfectly on its own and the
system still does not.}
